% latexmk (lualatex)
\documentclass[report,a4j,11pt]{ltjsbook}
\usepackage[luatex]{graphicx}
%\usepackage{graphicx}
\usepackage{url}
\usepackage{float}
\usepackage{subcaption}
\usepackage{amsmath}
\usepackage{comment}
%This file is encoded in UTF-8.


%\usepackage{txfonts}	% Times系フォントを利用するときはこのコメントを外すo
%\newcommand{\NiigataUnivMastersThesis}{}% !!!修論はコメントを外せ!!!

\和暦	% 日付を和暦で表示

%%% 均等割りコマンドを定義 %%%
\newcommand{\kintou}[2]{%
\leavevmode
\hbox to #1{
\ltjsetparameter{kanjiskip=0pt plus 1fill minus 1fill}
\ltjsetparameter{xkanjiskip=\ltjgetparameter{kanjiskip}}
#2}}

%%% 均等割りされて表示されるように設定 %%%
\renewcommand{\contentsname}{\kintou{5\zw}{目次}}
\renewcommand{\listfigurename}{\kintou{5\zw}{図目次}}
%\renewcommand{\listtablename}{\kintou{5\zw}{表目次}}
\renewcommand{\bibname}{\kintou{5\zw}{参考文献}}

\setcounter{tocdepth}{2}	% subsubsectionも目次に表示されるよう設定


%%% 【修論のみ】修論の表紙 %%%
%\ifdefined \NiigataUnivMastersThesis
%\title{卒論／修論テンプレート}
%\author{
%    \kintou{3\zw}{氏名}
%	\hspace{2\zw}
%	\kintou{6\zw}{新潟　太郎}
%	\\
%	\\
%	新潟大学 大学院自然科学研究科 博士前期課程 電気情報工学専攻
%}
%\date{\today}
%\fi

\begin{document}

\ifdefined \NiigataUnivMastersThesis
%\maketitle	% 【修論のみ】修論表紙をつける
\fi

%%% これより前付 %%%
\frontmatter
\ifdefined \NiigataUnivMastersThesis
%\chapter{\kintou{5\zw}{概要}}
修論の場合は，ここに概要 (Abstract) を書きます．
	% 【修論のみ】修論Abstractをつける
\fi
\chapter{\kintou{4\zw}{謝辞}}
はじめに,本研究を進めるにあたり,御指導頂いた村松正吾教授に深く感謝致します.村松教授には,日頃から的確かつ有益な熱い指導を頂きました.心より感謝申し上げます.
新潟大学災害・復興科学研究所の安田浩保教授には，本研究で扱った実験データをいただきましたことを心より感謝致します．
また，佐々木重信教授,金ミンソク准教授には多くのご意見を頂き感謝致します．
おわりに,これまで数々のご支援や助言を頂きました長山知司先輩，伊郷恭平先輩をはじめ,村松研究室の諸先輩方へ感謝の意を表します．本研究は，平成28年度 新潟大学U-goグラント（代表：早坂圭司）の助成による．
	% 謝辞
\tableofcontents	% 目次
\listoffigures		% 図目次
\listoftables		% 表目次

%%% これより本文 %%%
\mainmatter
\chapter{\kintou{5\zw}{はじめに}}

\section{背景}
近年，～

\section{問題}
流路変動の要因に～

\section{目的}
本研究では，～

\section{手段}
システムに線形性と～

\section{構成}
本論文の構成は以下の通りである．
第2章では，～
第3章では，～
第4章では，～
第5章では，～






	% 第1章：序論
\include{nisyou}		% 第2章：水面・河床同時計測の概要
\include{sansyou}   %第3章：水面・河床応答のシステム同定の提案
\include{yonsyou}   %第4章：性能評価
\include{gosyou}    %第5章：おわりに
%%% これより付録 %%%
\appendix
\chapter{最小二乗法}
この付録では，～

%%% 後付け %%%
\backmatter
\include{reference}	%参考文献

\end{document}
